% Part:sets-functions-relations
% Chapter: size-of-sets
% Section: pairing

\documentclass[../../../include/open-logic-section]{subfiles}

\begin{document}

\olfileid{sfr}{siz}{pai}
\olsection{جوڑی فنکشن تے کوڈ}

\begin{explain}
کانتور دا زگ زیگ طریقہ $\Nat^n$ دے گنتی وار فہرست بنن جوگ ہون نوں
اکھیں ویکھن جوگ بنا دیندا اے۔ پر آؤ اپنے $\Nat^2$ نوں وکھاؤن والے جدول
اُتے دھیان دیئیے۔ جدول وچ زگ زیگ لکیر دے نال نال چل کے تے مقام گن کے
اسی جانچ سکدے آں کہ $\tuple{1,2}$ نوں عدد~$7$ نال جوڑیا گیا اے۔ پر جے
اسی ایہ گل ہور سدھے ڈھنگ نال حساب کر سکدے، تاں چنگا ہوندا۔ یعنی ساڈے
کول زگ زیگ گنتی وار فہرست دا \emph{اُلٹ فنکشن}،
$g\colon \Nat^2 \to \Nat$، ہونا چنگا ہوندا، ایہو جہا کہ
\[
g(\tuple{0,0}) = 0, \;
g(\tuple{0,1}) = 1, \;
g(\tuple{1,0}) = 2, \; \dots,
g(\tuple{1,2}) = 7, \; \dots
\]
ایس نال اسی ٹھیک ٹھیک حساب کر سکدے آں کہ $\tuple{n, m}$ ساڈی گنتی وار
فہرست وچ کتھے آوے گی۔

اصل وچ اسی دو گلاں اُتے غور کر کے $g$ نوں سدھے ڈھنگ نال تعریف کر سکدے
آں۔ پہلی گل: جے $n$ویں قطار تے $m$ویں کالم وچ قدر~$v$ اے، تاں
$(n+1)$ویں قطار تے $(m-1)$ویں کالم وچ قدر $v + 1$ اے۔ دوجی گل:
ساڈی گنتی وار فہرست دی پہلی قطار تکونی عدداں اُتے مشتمل اے، جیہڑی
$0$، $1$، $3$، $6$ وغیرہ توں شروع ہوندی اے۔ $k$واں تکونی عدد،
$< k$ قدرتی عدداں دا جوڑ اے، جیہڑا $k(k+1)/2$ دے طور تے حساب کیتا جا
سکدا اے۔ ایہ دوویں گلاں کٹھیاں کر کے اگلے فنکشن اُتے غور کرو:
\[
  g(n,m) = \frac{(n+m+1)(n+m)}{2} + n
\]
اسی اکثر صرف $g(n, m)$ لکھدے آں، $g(\tuple{n, m})$ دی بجائے، کیوں جے
ایہ اکھاں نوں سوکھا لگدا اے۔ ایہ تہانوں آکھدا اے کہ پہلاں
$(n+m)^\text{th}$ تکونی عدد کڈھو، تے فیر اوہدے وچ $n$ جوڑ دیو۔ ایہ
جدول نوں ٹھیک اوسے ڈھنگ نال بھردا اے جیویں اسی چاہندے آں۔ سو خاص طور
تے جوڑی $\tuple{1, 2}$ نوں $\frac{4 \times 3}{2} +
1 = 7$ ول بھیجیا جاندا اے۔

ایہ فنکشن $g$، جوڑیاں دے سیٹ دی اک گنتی وار فہرست دا
\emph{اُلٹ فنکشن} اے۔ ایہو جہے فنکشناں نوں \emph{جوڑی فنکشن} آکھیا
جاندا اے۔
\end{explain}

\begin{defn}[جوڑی فنکشن]
  فنکشن $f\colon A \times B \to \Nat$ اک حسابی
  \emph{جوڑی فنکشن} اے جے $f$، اک بہ اک ہووے۔ اسی ایہ وی آکھدے آں
  کہ $f$، $A \times B$ نوں \emph{کوڈ کردا اے}، تے $f(x,y)$،
  $\tuple{x,y}$ دا \emph{کوڈ} اے۔
\end{defn}

\begin{explain}
اسی جوڑی فنکشناں نوں، مثلاً، قدرتی عدداں دیاں جوڑیاں کوڈ کرن لئی ورت
سکدے آں؛ یا دوجے لفظاں وچ، اسی عنصراں دی ہر \emph{جوڑی} نوں اک
\emph{اکلے} عدد نال پیش کر سکدے آں۔ جوڑی فنکشن دا اُلٹ فنکشن ورت کے
اسی عدد نوں \emph{ڈی کوڈ} کر سکدے آں، یعنی پتا لا سکدے آں کہ اوہ کیہڑی
جوڑی نوں پیش کردا اے۔
\end{explain}

\begin{prob}
سارے غیر منفی ناطق عدداں دے سیٹ دی اک گنتی وار فہرست دیو۔
\end{prob}

\begin{prob}
وکھاؤ کہ $\Rat$، !!{enumerable} اے۔ یاد رکھو کہ ہر ناطق عدد نوں
$z/m$ کسر دے طور تے لکھیا جا سکدا اے، جتھے $z \in \Int$ تے
$m \in \Nat^+$۔
\end{prob}

\begin{prob}
$\Bin^*$ دی اک گنتی وار فہرست تعریف کرو۔
\end{prob}

\begin{prob}
اپنے ابتدائی منطق کورس توں یاد کرو کہ ہر ممکن صداقتی جدول اک صداقتی
فنکشن نوں ظاہر کردا اے۔ دوجے لفظاں وچ، صداقتی فنکشن
$\Bin^k \to \Bin$ والے سارے فنکشن نیں، کسے~$k$ لئی۔ ثابت کرو کہ سارے صداقتی فنکشناں
دا سیٹ گنتی وار فہرست بنن جوگ اے۔
\end{prob}

\begin{prob}
وکھاؤ کہ کسے وی لامتناہی !!{enumerable} سیٹ دے سارے متناہی ذیلی سیٹاں
دا سیٹ !!{enumerable} اے۔
\end{prob}

\begin{prob}
$\Nat$ دے کسے ذیلی سیٹ نوں \emph{ہم متناہی} عین اوس ویلے آکھیا جاندا اے
جدوں اوہ قدرتی عدداں دے کسے متناہی ذیلی سیٹ دا $\Nat$ دے اندر متمم
ہووے؛ یعنی $A \subseteq \Nat$ عین اوس ویلے ہم متناہی اے جدوں
$\Nat\setminus A$ متناہی ہووے۔
% OLSIZ-002 ماخذ دی درستی: متمم قدرتی عدداں دے سیٹ دے اندر کسے متناہی ذیلی سیٹ دا اے؛ انگریزی عبارت وچ ایہ نسبت رہ گئی سی۔
$I$ اوہ سیٹ ہووے جیہدے !!{element}s ٹھیک $\Nat$ دے متناہی تے ہم متناہی
ذیلی سیٹ نیں۔ وکھاؤ کہ $I$، !!{enumerable} اے۔
\end{prob}

\begin{prob}
وکھاؤ کہ !!{enumerable} سیٹاں دا !!{enumerable} اتحاد وی
!!{enumerable} اے۔ یعنی جدوں وی $A_1$، $A_2$، \dots{} سیٹ ہون، تے
ہر $A_i$، !!{enumerable} ہووے، تاں اوہناں ساریاں دا اتحاد
$\bigcup_{i=1}^\infty A_i$ وی !!{enumerable} اے۔ [نوٹ: ایہ اوکھا اے!]
\end{prob}

\begin{prob}
$f \colon A \times B \to \Nat$ کوئی وی جوڑی فنکشن ہووے۔ وکھاؤ کہ
$f$ دا اُلٹ فنکشن $A \times B$ دی اک گنتی وار فہرست اے۔
\end{prob}

\begin{prob}
ایہو جہا فنکشن صاف صاف دسو جیہڑا $\Nat^3$ نوں کوڈ کردا ہووے۔
\end{prob}

\end{document}
